\problem{P02}{Antithetic Estimation under Symmetry}{Symmetry and variance}
Suppose $S$ and $-S$ have the same distribution, and write $\tau=\Prob(S=0)$. To estimate $p=\Prob(S>0)$, compare
\[
A=\frac{\ind\{S_1>0\}+\ind\{S_2>0\}}2,
\qquad B=\frac{\ind\{S>0\}+\ind\{-S>0\}}2,
\]
where $S_1,S_2$ are independent copies of $S$. Prove that both estimators are unbiased, find their variances, and determine when the antithetic estimator has no random error.
\solution Symmetry gives $p=\Prob(S<0)=(1-\tau)/2$, so $\E A=\E B=p$. Independence of the two indicators in $A$ yields
\[
\Var A=\frac{p(1-p)}2=\frac{1-\tau^2}{8}.
\]
The variables $S$ and $-S$ cannot both be positive. Unless $S=0$, exactly one is positive. Consequently,
\[
B=\frac{1-\ind\{S=0\}}2,
\qquad \Var B=\frac{\tau(1-\tau)}4,
\]
and hence
\[
\boxed{\Var A-\Var B=\frac{(1-\tau)^2}{8}\ge0.}
\]
For $0\le\tau<1$, the variance ratio is $2\tau/(1+\tau)$. In particular, $\tau=0$ implies $B\equiv1/2$: pairing removes all random error. If $\tau=1$, both estimators are identically zero, so their variance ratio is undefined.
\insight The improvement comes from negative dependence. Two evaluations with the same marginal law need not be independent; pairing opposite signs leaves only the uncertainty of the atom at zero.
